\section{A weighted presentation for the pinched fiber}
\label{sec:involution}

We now equip the free pro-\(2\) group \(L\) with a filtration adapted to
the covering shifts.  In this filtration the involution relations from
the preceding section can be imposed with a sufficiently small contribution to the weighted
relation sum.  We also show that the degree bounds on \(F\) survive the map
\(\Phi:F\to L\rtimes H\). Although the pinched sphere group is
already free of rank $24$, an arbitrary basis of that quotient would
not retain these degree bounds. Keeping $48$ generators and imposing
$24$ carefully chosen relations preserves the filtration needed to
control all later monodromy relations.

\subsection{A basis adapted to the two translations}

Recall that \(L\) is free on \(a_{h,j}\), where
\(h\in H=(\Z/4)^2\) and \(j\in\{0,1,2\}\).  The commuting shifts
\(T_1,T_2\) translate the two coordinates of \(h\).  The involution is
\[
 \sigma(a_{h,j})=a_{b_j-h,j}^{-1},
 \qquad b_j\bmod 2\ne0.
\]
Thus \(\sigma T_r\sigma=T_r^{-1}\) for \(r=1,2\).
For an automorphism \(T\), write
\(\delta_T(g)=T(g)g^{-1}\), and define
\begin{equation}\label{eq:shift-basis}
 u_{i\ell,j}=\delta_{T_1}^{\,i}\delta_{T_2}^{\,\ell}(a_{0,j}),
 \qquad 0\le i,\ell<4.
\end{equation}

\begin{lemma}\label{lem:shift-basis}
The elements in \eqref{eq:shift-basis} form a free pro-\(2\) basis of
\(L\).  Give \(u_{i\ell,j}\) weight \(i+\ell+1\), and denote the
resulting Magnus valuation by \(\nu_L\).  Each of
\(T_1-\id\), \(T_2-\id\), and \(\sigma-\id\) raises the degree of the
completed Magnus algebra by at least one.
\end{lemma}

\begin{proof}
In the Frattini quotient, the span of the generators in one block is
the permutation module
\[
 R=\F_2[H]
   =\F_2[\xi,\eta]/(\xi^4,\eta^4),
 \qquad T_1=1+\xi,\quad T_2=1+\eta.
\]
The images of \(u_{i\ell,j}\) are its monomial basis
\(\xi^i\eta^\ell\).  The Frattini basis criterion therefore gives the
first assertion.  We henceforth use the stated weights.

The proof has three steps: control the last difference of a cyclic
shift, propagate the two translation bounds, and then treat sign.
We first establish the endpoint estimate needed for a cyclic shift.
Suppose \(T^4=\id\), put \(w_i=\delta_T^i(w_0)\), and assume
\begin{equation}\label{eq:cyclic-hypothesis}
 \nu_L(w_i)\ge d+i\qquad(0\le i\le3),
 \qquad d\ge1.
\end{equation}
Then
\begin{equation}\label{eq:cyclic-endpoint}
 \nu_L(w_4)\ge d+4.
\end{equation}
This assertion does not assume that \(T\) preserves the filtration.
To prove it, take formal free generators \(z_0,z_1,z_2,z_3\) of weights
\(d,d+1,d+2,d+3\), and define the triangular automorphism
\[
 T'(z_i)=z_{i+1}z_i\quad(0\le i<3),
 \qquad T'(z_3)=z_3.
\]
By \Cref{lem:substitution}, \(T'-\id\) raises algebra degree by one.
In characteristic two,
\[
 (T')^4-\id=(T'-\id)^4,
\]
so \(T'^4(z_0)z_0^{-1}\) has valuation at least \(d+4\).
Evaluation at \(z_i=w_i\) preserves this bound by
\eqref{eq:cyclic-hypothesis}.  The positive word \(T'^3(z_0)\) contains
\(z_3\) exactly once, as its first letter.  In passing from the
evaluated formal fourth iterate to the actual fourth iterate, the only
changed substitution is
\[
 T(w_3)=w_4w_3
 \quad\hbox{in place of}\quad
 T'(z_3)=z_3.
\]
Consequently \(T^4(w_0)=w_4\,T'^4(z_0)|_{z_i=w_i}\).
Since \(T^4(w_0)=w_0\), the bound for the evaluated formal iterate
proves \eqref{eq:cyclic-endpoint}.

For each fixed \(\ell,j\), apply this estimate to
\(w_i=u_{i\ell,j}\), \(T=T_1\), and \(d=\ell+1\).
For \(i<3\), the required difference is exactly
\(\delta_{T_1}(u_{i\ell,j})=u_{i+1,\ell,j}\);
\eqref{eq:cyclic-endpoint} supplies the case \(i=3\).
Thus \(T_1-\id\) raises degree by one on every basis element, and
therefore on the algebra.

The same argument gives the \(T_2\) bounds on the initial sequence
\(u_{0\ell,j}\).  They propagate along the \(T_1\)-sequences as follows.
If \(T_2(z)=vz\), commutativity of the shifts gives
\begin{equation}\label{eq:commuting-difference}
 T_2(\delta_{T_1}z)
   =T_1(v)(\delta_{T_1}z)v^{-1}.
\end{equation}
Suppose \(z\) has assigned weight \(d\), the element \(v\) has
valuation at least \(d+1\), and \(\delta_{T_1}z\) has valuation at
least \(d+1\).  The already proved \(T_1\) bound gives
\(\nu_L(T_1(v)v^{-1})\ge d+2\).  Conjugating
\(\delta_{T_1}z\) by \(v\) changes it only in degree at least
\(2d+2\ge d+2\).  Equation \eqref{eq:commuting-difference} therefore
gives the next \(T_2\) bound.  Induction proves that \(T_2-\id\)
raises degree by one on the full basis.

It remains to treat \(\sigma\).  Powers of either shift preserve the
filtration and differ from the identity by an operator raising degree
by one.  Since
\(\sigma(a_{0,j})=T_1^{(b_j)_1}T_2^{(b_j)_2}(a_{0,j})^{-1}\),
and inversion has the same linear term in characteristic two, we have
\[
 \nu_L\bigl(\sigma(a_{0,j})a_{0,j}^{-1}\bigr)\ge2.
\]
We propagate this estimate along the two difference sequences.
For either shift \(T\), its inverse preserves the filtration, and
\[
 T^{-1}-T=T^{-1}(\id-T^2),
 \qquad T^2-\id=(T-\id)^2.
\]
Hence, if \(\nu_L(z)\ge d\), the elements
\(\delta_{T^{-1}}z\) and \(\delta_Tz\) agree modulo degree \(d+2\).
If \(\sigma(z)=vz\), the identity
\(\sigma T=T^{-1}\sigma\) gives
\begin{equation}\label{eq:inverting-difference}
 \sigma(\delta_Tz)
   =T^{-1}(v)(\delta_{T^{-1}}z)v^{-1}.
\end{equation}
When \(\nu_L(v)\ge d+1\), the inverse-shift bound and the
commutator estimate show that the right side differs from
\(\delta_{T^{-1}}z\) only in degree at least \(d+2\).
The preceding comparison then replaces
\(\delta_{T^{-1}}z\) by \(\delta_Tz\).
Starting with \(a_{0,j}\), first apply this argument along
\(\delta_{T_2}\) and then along \(\delta_{T_1}\).
It proves that \(\sigma\) changes each basis element of weight \(w\)
only in degree at least \(w+1\).  A final application of
\Cref{lem:substitution} proves the assertion on the completed algebra.
\end{proof}

\subsection{Twenty-four relations impose the involution}

There are 24 pairs of generators exchanged, up to inversion, by sign.
Choosing one relation per pair in the original basis would impose the
involution, but would lose the higher degrees gained from the shifts.
We must choose 24 relations whose degrees are measured in the new
basis. The Frattini quotient will identify those relations, and its
generation criterion will show that they impose the full group action.

\begin{proposition}\label{prop:involution}
Let \(L\) be the free pro-\(2\) group on \(a_{h,j}\), with
\(h\in(\Z/4)^2\) and \(0\le j\le2\), and let
\(\sigma(a_{h,j})=a_{b_j-h,j}^{-1}\), where \(b_j\bmod2\ne0\).
Equip \(L\) with the basis and weights of \Cref{lem:shift-basis}.
There are 24 elements \(g\in L\), with assigned positive weights
\(w_g\), such that
\[
 \nu_L(g)\ge w_g,\qquad
 \nu_L(r_g)\ge w_g+1,
 \qquad r_g=\sigma(g)g^{-1},
\]
and the closed normal subgroup generated by the \(r_g\) imposes
\(\sigma=\id\) on all of \(L\).  The basis polynomial and the
polynomial of the chosen weights are
\begin{equation}\label{eq:weight-polynomials}
 \begin{split}
 P(t)&=\sum_{i,\ell,j}t^{i+\ell+1}
       =3t(1+t+t^2+t^3)^2,\\
 G(t)&=\sum_g t^{w_g}
       =3t(1+t)(1+t^2)^2.
 \end{split}
\end{equation}
In particular, \(P=(1+t)G\), and these involution relations have
degree cost at most \(tG(t)\).
\end{proposition}

\begin{proof}
The full Magnus degree bounds are already proved. We now use the
Frattini quotient to choose enough relations to impose the whole
involution; generation is detected in this linear quotient.
Consider one block of the Frattini quotient, identified with
\(R=\F_2[\xi,\eta]/(\xi^4,\eta^4)\) as above.
Give a homogeneous polynomial of degree \(n\) weight \(n+1\).
Write \(b=(b_1,b_2)\) for the shift of this block.
The action of \(\sigma\) on \(R\) is
\begin{equation}\label{eq:frattini-involution}
 f(\xi,\eta)\longmapsto
 (1+\xi)^{b_1}(1+\eta)^{b_2}
 f\bigl((1+\xi)^{-1}-1,(1+\eta)^{-1}-1\bigr).
\end{equation}
Here inversion of a group generator disappears in the mod-\(2\)
Frattini quotient.  Since
\((1+\xi)^{-1}-1=\xi+\xi^2+\cdots\), expansion of
\eqref{eq:frattini-involution} shows that the part of
\(\sigma-\id\) raising weight by exactly one is
\begin{equation}\label{eq:frattini-differential}
 d=(\bar b_1\xi+\bar b_2\eta)
       +\xi^2\frac{\partial}{\partial\xi}
       +\eta^2\frac{\partial}{\partial\eta},
 \qquad (\bar b_1,\bar b_2)=b\bmod2.
\end{equation}
The first term means multiplication.

Set \(u=\bar b_1\xi+\bar b_2\eta\), which is nonzero, and choose a
linear coordinate \(v\) independent of \(u\) over \(\F_2\).
The fourth-power relations become \(u^4=v^4=0\).
Moreover, the vector field
\(\xi^2\partial_\xi+\eta^2\partial_\eta\) sends any linear form
over \(\F_2\) to its square.  In these coordinates,
\begin{equation}\label{eq:even-coordinate-differential}
 R=\F_2[u,v]/(u^4,v^4),
 \qquad d=u+u^2\frac{\partial}{\partial u}
                +v^2\frac{\partial}{\partial v}.
\end{equation}
The \(u\)-part sends \(1\) to \(u\), \(u^2\) to \(u^3\), and
odd powers to zero.  Its square is zero, as is the square of the
\(v\)-part.  The two parts commute, so \(d^2=0\).

Take the homogeneous subspace
\begin{equation}\label{eq:even-complement}
 C=\operatorname{span}_{\F_2}
       \{u^a v^b:a\in\{0,2\},\ 0\le b<4\}.
\end{equation}
On \(C\), the odd-\(u\) component of \(d\) is multiplication by
\(u\), an isomorphism from \(C\) to the span of the odd-\(u\)
monomials.  Thus \(d|_C\) is injective and
\begin{equation}\label{eq:acyclic-frattini}
 R=C\oplus dC,\qquad \ker d=\im d=dC.
\end{equation}
Indeed the direct-sum assertion follows by projecting to the
odd-\(u\) monomials; both summands have dimension eight.
It then follows from \(d^2=0\) that the image and kernel, each of
dimension eight, coincide.

For each of the eight monomials in \eqref{eq:even-complement},
expand it in the original homogeneous monomials \(\xi^i\eta^\ell\).
Choose \(g\) to be a product of the corresponding
\(u_{i\ell,j}\), each appearing with its coefficient in \(\F_2\),
in any fixed order.  Its assigned weight \(w_g\) is the polynomial
degree plus one.  Then \(\nu_L(g)\ge w_g\), and
\Cref{lem:shift-basis} gives \(\nu_L(r_g)\ge w_g+1\).
In the graded Frattini quotient the leading vectors of \(g,r_g\)
are respectively the chosen basis of \(C\) and its image under
\(d\).  By \eqref{eq:acyclic-frattini} these form a homogeneous
basis.  The actual Frattini vectors therefore form a basis as well:
any linear dependence would give one between their lowest nonzero
weight components.  Doing this in all three blocks shows that the
48 elements \(g,r_g\), equivalently \(g,\sigma(g)\), generate \(L\).

Let \(N\) be the closed normal subgroup generated by these 24
elements \(r_g\).  It is invariant under \(\sigma\), because
\[
 \sigma(r_g)=g\sigma(g)^{-1}=r_g^{-1}.
\]
The quotient \(L/N\) is generated by the images of the \(g\), since
there \(\sigma(g)=g\); its induced involution fixes those generators
and hence is the identity.  Thus \(N\) is exactly the closed normal
subgroup imposing \(\sigma=\id\).

Finally, in each block the chosen monomials have weight polynomial
\[
 t(1+t^2)(1+t+t^2+t^3)=t(1+t)(1+t^2)^2.
\]
Adding the three blocks gives \eqref{eq:weight-polynomials}.
Since each \(r_g\) has degree at least \(w_g+1\), its contribution
to the relation bound is at most \(t^{w_g+1}\) for \(0<t<1\).
\end{proof}

\subsection{Transferring the degree bounds}
\label{sec:transfer}

The involution relations leave the positive polynomial \(G=P-tG\)
in the generator-minus-relation count.
To control the remaining monodromy relations, we need to
compare the two filtrations.

Write $\widehat K$ for the closure of $K$ in $\widehat F$.
The normal subgroup $K\triangleleft F$ has the induced pro-2 topology: if $N\triangleleft K$ has 2-power index, its $F$-core $N_0$ is the intersection of finitely many conjugates of $N$. The group $K/N_0$ embeds in a product of finite 2-groups, and $F/N_0$ is an extension of $F/K$ by $K/N_0$, hence is also a finite 2-group. This proves the required cofinality and identifies the pro-2 completion of $K$ with its closure in $\widehat F$.

\begin{lemma}\label{lem:transfer}
Let \(\Phi:\widehat F\to L\rtimes H\) be the continuous extension of
the homomorphism from the preceding section:
\[
 x\longmapsto T_1,\qquad y\longmapsto T_2,\qquad
 a_j\longmapsto a_{0,j},\qquad c_0,c_1\longmapsto1.
\]
If \(f\in\widehat F\) has \(\nu_F(f)\ge n\), its \(L\)-component
under \(\Phi\) has \(\nu_L\ge n\).  In particular, this holds for
the homomorphism \(\Phi|_{\widehat K}:\widehat K\to L\).
\end{lemma}

\begin{proof}
Let \(A_L\) be the completed Magnus algebra of \(L\).
Represent \(L\rtimes H\) on \(A_L\) by letting \(l\in L\) act by
left multiplication and letting \(H\) act by its shifts.
By \Cref{lem:shift-basis}, the operators representing \(x-1\) and
\(y-1\) raise degree by one.  So does the operator representing
\(a_j-1\), which is left multiplication by \(a_{0,j}-1\).
The operators representing \(c_0-1,c_1-1\) are zero, and in
particular satisfy their required weight-two bounds.

Substitute these operators into the weighted Magnus algebra of
\(\widehat F\).  This substitution converges: a monomial of weight
at least \(n\) becomes an operator raising degree by at least \(n\).
Consequently the operator representing \(f-1\) raises degree by
at least \(n\).  If \(\Phi(f)=(l,h)\), its action on \(1\in A_L\)
is \(l\), since every shift fixes \(1\).  Evaluating the operator
for \(f-1\) on \(1\) therefore gives \(l-1\), proving
\(\nu_L(l)\ge n\).  On \(\widehat K\) the \(H\)-component is
trivial, which proves the last assertion.
\end{proof}
